\subsection{Explanation layer}\label{sec:xai}
\subsubsection{Stability and certification}
Each alert carries a SHAP attribution \cite{lundberg2017}, rendered in plain language and split into evidence for and against the verdict. Under a 1\,\% Gaussian perturbation, SHAP's top-5 features have a Jaccard stability of 0.83 (Spearman $+0.92$; Review~2 quoted 0.87 from an earlier run of the same evaluation, and we report the value the current evaluation script produces), and LIME's \cite{ribeiro2016} 0.27 (Spearman $+0.21$), while the model's own prediction flips in 0\,\% of cases, so the instability belongs to the explainer. Raising LIME's sample count from 1,000 to 12,000 moves its Jaccard only within 0.16 to 0.38 (0.16, 0.38 and 0.33, and its apparent fidelity falls from $+0.349$ to $+0.023$ as sampling converges. LIME was demoted to an internal cross-check. Each SHAP explanation is certified before display by ablating its top-cited features and checking that the decision moves more than for random features; 82\,\% of 40 explanations certify, and the rest are shown labelled unverified rather than hidden.

\subsubsection{Human evaluation of explanation stability}
Certification only helps if a reader notices it. We designed a within-participant reader study with pre-registered hypotheses: (H1) readers state the reason for a certified alert correctly more often than for an uncertified one; (H2) readers trust uncertified alerts less; (H3) contradictory alerts, whose top-ranked evidence argues against the verdict, lead readers to cite counter-evidence as the reason. Ten to fifteen participants without exposure to the model each read 20 real alerts produced by the trained ensemble (8 certified, 6 uncertified, 6 contradictory) in a seeded per-participant order, and for each state the verdict, the system's confidence band, the reason (a checklist of the cited measurements plus two distractors that the alert never cites, and one free sentence), the action they would take, their own confidence and their trust. Time to decision is recorded. A reason is scored correct when the reader selects at least one measurement the alert cites for its verdict and no counter-evidence or distractor, the pass rule fixed before the study. Analysis compares participant-level means per category with paired Wilcoxon tests, Bonferroni-corrected, and actions by category with a chi-square test; free-text reports of unreliable alerts are coded independently by two team members with Cohen's kappa.

\pending{Run the protocol (reader\_study/PROTOCOL.md; study app; 25 minutes per participant), export the CSV, run reader\_study/analyse\_reader\_study.py and insert participant demographics, accuracy, confidence, time and trust per category here. No result is reported until the sessions have been run.}
